\subsection{\problemblock{*Increment} Data Block}\label{sec:block-increment}

Normally the values for a vehicle's state and time are continuous across trajectory
segments. The state or time can, however, be changed instantaneously to form a
discontinuity by including an \problemblock{*increment} data block. This is used
most often for weight changes, such as dropping a weapon or staging a missile. For
example,

\begin{lstlisting}[style=taosmanual]
*increment wt=-325
\end{lstlisting}

instantaneously reduces the weight by 325~lb when default units are used.

The \problemblock{*increment} block is identical to the
\problemblock{*reset} block described in \cref{sec:block-reset}, except that
\problemblock{*reset} sets a time or state variable equal to an input value,
whereas \problemblock{*increment} adds the input value to the current value. In
other words, \problemblock{*increment} applies a positive or negative delta to
the current time or state variable.

If both \problemblock{*reset} and \problemblock{*increment} blocks are
present in a segment, values are reset first and then incremented. Multiple reset and
increment blocks can be given; they are executed in the order in which they appear.

The variables that can be changed are listed in \cref{tab:increment-variables}.
Position and velocity are maintained in several coordinate systems, so the variables
are grouped into consistent sets. Position or velocity values in one increment block
must belong to the same coordinate system. For example, if
\statevar{xecicdt} is incremented, then \statevar{yecicdt} and
\statevar{zecicdt} may also be changed, but velocity in another representation,
such as \statevar{vel}, cannot be changed in the same block. This restriction
applies to every position and velocity set in the table.

\begin{longtable}{@{}>{\ttfamily}lp{0.68\linewidth}@{}}
  \caption{Variables accepted by \texttt{*increment} and \texttt{*reset}.}
  \label{tab:increment-variables}\\
  \toprule
  \normalfont Variable & Description \\
  \midrule
  \endfirsthead
  \toprule
  \normalfont Variable & Description \\
  \midrule
  \endhead
  alt       & Altitude \\
  long      & Longitude (with geodetic latitude and altitude) \\
  latgd     & Geodetic latitude \\
  rcm       & Distance from earth center \\
  long      & Longitude (with geocentric radius and latitude) \\
  latgc     & Geocentric latitude \\
  xecfc     & Earth-centered, earth-fixed $x$ coordinate \\
  yecfc     & Earth-centered, earth-fixed $y$ coordinate \\
  zecfc     & Earth-centered, earth-fixed $z$ coordinate \\
  xecic     & Earth-centered inertial $x$ coordinate \\
  yecic     & Earth-centered inertial $y$ coordinate \\
  zecic     & Earth-centered inertial $z$ coordinate \\
  dxb       & Position increment along body $x$ axis \\
  dyb       & Position increment along body $y$ axis \\
  dzb       & Position increment along body $z$ axis \\
  vel       & Velocity \\
  gamgc     & Geocentric flight-path angle \\
  psigc     & Geocentric heading angle \\
  vel       & Velocity \\
  gamgd     & Geodetic flight-path angle \\
  psigd     & Geodetic heading angle \\
  xecfcdt   & Earth-centered, earth-fixed velocity $x$ component \\
  yecfcdt   & Earth-centered, earth-fixed velocity $y$ component \\
  zecfcdt   & Earth-centered, earth-fixed velocity $z$ component \\
  xecicdt   & Earth-centered inertial velocity $x$ component \\
  yecicdt   & Earth-centered inertial velocity $y$ component \\
  zecicdt   & Earth-centered inertial velocity $z$ component \\
  wt        & Weight \\
  mass      & Mass \\
  fuel      & Fuel used \\
  time      & Time \\
  tseg      & Segment time \\
  tmark     & Mark time \\
  range     & Ground range \\
  grseg     & Segment ground range \\
  grmark    & Mark ground range \\
  plength   & Path length \\
  plseg     & Segment path length \\
  plmark    & Mark path length \\
  iip\_beta & Ballistic coefficient for initial-impact-point calculations \\
  \bottomrule
\end{longtable}

Weight and mass are similar. If one is given, the other cannot be given because both
refer to the same state variable. Fuel, on the other hand, may be reset or incremented
at any time.

The time, range, and path-length variables are independent of one another, so any
combination is valid. The segment variables \statevar{tseg}, \statevar{grseg}, and
\statevar{plseg} are automatically reset to zero at the beginning of a segment. The
increment values are applied after that automatic reset.

A problem with one trajectory can contain a discontinuity in absolute time, so
\statevar{time} may be reset or incremented. Problems with more than one
trajectory use time to relate the vehicles to one another; a time discontinuity can
therefore cause errors. The variable \statevar{time} should not be reset or
incremented in a multiple-trajectory problem. The relative-time variables
\statevar{tseg} and \statevar{tmark} do not have this problem.

The mark variables \statevar{tmark}, \statevar{grmark}, and
\statevar{plmark} are useful for measuring time, range, and path length from an
intermediate point. At the beginning of a trajectory they correspond to
\statevar{time}, \statevar{range}, and \statevar{plength}. They can be restarted
at any segment with \problemblock{*reset} or changed with
\problemblock{*increment}.

The ballistic coefficient for the initial-impact-point calculation may also be reset or
incremented at any time. It is used only by the IIP calculation.

\taossubhead{Object Deployment}

The variables \statevar{dxb}, \statevar{dyb}, and \statevar{dzb} are used when an
object is deployed or a weapon is dropped and one trajectory is initialized from
another. The deployed object's initial position is usually offset slightly from the
vehicle center of mass. The three variables give the object's position relative to the
body coordinate system before deployment.

A deployed object is generally initialized from the beginning of an existing segment.
For example,

\begin{lstlisting}[style=taosmanual]
*trajectory 2 Obj2 start on 1

  *initial from trajectory 1, segment 6

  *segment 1 deployment
    *reset wt=50
    *increment velibx=5
    *increment dxb=0.5 dyb=3.2 dzb=0.35
    *fly pitchi=*
    *fly yawi=*
    *fly rolli=*
    *when tseg>0 goto 2
\end{lstlisting}

initializes trajectory 2 from the beginning of segment 6 in trajectory 1 and
represents an object deployment in space. The initial state is copied from trajectory
1 and can be adjusted with reset and increment blocks in the first segment.

In this example the new object's weight is set to 50~lb, and it is deployed with a
5~ft/s velocity increment along the body $x$ axis. Its position is adjusted with
\statevar{dxb}, \statevar{dyb}, and \statevar{dzb} because its center of mass is
slightly offset from the original vehicle. These three variables always act as position
increments, so they have the same effect in a reset or increment block.

The guidance rules hold a constant body attitude in inertial space. This ensures that
the object's body attitude matches the original vehicle's attitude, so the body-axis
position and velocity increments are applied in the correct directions. The same
result can be obtained by using the original vehicle's guidance rules in the first
segment. A final condition of \texttt{tseg>0} ends the deployment segment
immediately so that other guidance rules can be used for the remaining trajectory.
